\chapter*{Resumen}
¿Cuántas veces nos hemos enterado de un concierto, una exposición o una fiesta de barrio cuando ya había pasado? A veces por un cartel que vimos tarde, a veces porque alguien lo subió a Instagram y no preguntamos a tiempo, a veces simplemente porque nadie nos lo contó. Sevilla tiene una vida cultural y festiva extraordinaria, pero su oferta está dispersa: cada evento vive en su propia red social, en su propia web o en el boca a boca de quien ya lo conocía.

\textit{Sevillaneando} es una aplicación móvil que reúne toda esa oferta en un solo lugar, pero con una diferencia fundamental respecto a un simple directorio de eventos: está pensada para que la propia comunidad la construya y la haga crecer. Cualquier usuario puede publicar su propio evento ---un taller, una quedada, una ruta por el centro--- y compartirlo con quienes le siguen, o bien crearlo como evento privado con acceso restringido para que solo llegue a quien él decida. Los asistentes pueden chatear entre ellos antes y durante el evento, subir fotos en el momento, dejar reseñas y descubrir a otras personas con sus mismos intereses. La aplicación también sugiere planes personalizados según lo que a cada usuario le gusta y dónde se encuentra.

El resultado es una plataforma donde descubrir Sevilla deja de depender de la suerte o de a quién conozcas.

\vspace{.5cm}

\textbf{Palabras clave:} aplicación móvil, eventos culturales, Sevilla, ocio, comunidad, tiempo real, geolocalización, recomendaciones personalizadas
